% !TeX program = XeLaTeX
% !TeX root = ../../shloka.tex
% --meta--
% deity: Shiva
% composer: Adi Shankaracharya
% wiki_title: Mrityunjaya Manasa Puja Stotram
% --end-meta--

\sect{मृत्युञ्जयमानसपूजास्तोत्रम्}

\fourlineindentedshloka
{कैलासे कमनीयरत्नखचिते कल्पद्रुमूले स्थितं}
{कर्पूरस्फटिकेन्दुसुन्दरतनुं कात्यायनीसेवितम्}
{गङ्गातुङ्गतरङ्गरञ्जितजटाभारं कृपासागरं}
{कण्ठालङ्कृतशेषभूषणममुं मृत्युञ्जयं भावये}% ॥ १ ॥

\fourlineindentedshloka
{आगत्य मृत्युञ्जय चन्द्रमौले}
{व्याघ्राजिनालङ्कृत शूलपाणे}
{स्वभक्तसंरक्षणकामधेनो}
{प्रसीद विश्वेश्वर पार्वतीश}% ॥ २ ॥

\fourlineindentedshloka
{भास्वन्मौक्तिकतोरणे मरकतस्तम्भायुतालङ्कृते}
{सौधे धूपसुवासिते मणिमये माणिक्यदीपाञ्चिते}
{ब्रह्मेन्द्रामरयोगिपुङ्गवगणैर्युक्ते च कल्पद्रुमैः}
{श्रीमृत्युञ्जय सुस्थिरो भव विभो माणिक्यसिंहासने}% ॥ ३ ॥

\fourlineindentedshloka
{मन्दारमल्लीकरवीरमाधवी-}
{पुन्नागनीलोत्पलचम्पकान्वितैः}
{कर्पूरपाटीरसुवासितैर्जलैः}
{आधत्स्व मृत्युञ्जय पाद्यमुत्तमम्}% ॥ ४ ॥

\fourlineindentedshloka
{सुगन्धपुष्पप्रकरैः सुवासितैः}
{वियन्नदीशीतलवारिभिः शुभैः}
{त्रिलोकनाथार्तिहरार्घ्यमादरात्}
{गृहाण मृत्युञ्जय सर्ववन्दित}% ॥ ५ ॥

\twolineshloka
{हिमाम्बुवासितैस्तोयैः शीतलैरतिपावनैः}
{मृत्युञ्जय महादेव शुद्धाचमनमाचर}% ॥ ६ ॥

\fourlineindentedshloka
{गुडदधिसहितं मधुप्रकीर्णं}
{सुघृतसमन्वितधेनुदुग्धयुक्तम्}
{शुभकर मधुपर्कमाहर त्वं}
{त्रिनयन मृत्युहर त्रिलोकवन्द्य}% ॥ ७ ॥

\twolineshloka
{पञ्चास्त्रशान्त पञ्चास्य पञ्चपातकसंहर}
{पञ्चामृतस्नानमिदं कुरु मृत्युञ्जय प्रभो}% ॥ ८ ॥

\fourlineindentedshloka
{जगत्त्रयीख्यात समस्ततीर्थ-}
{समाहृतैः कल्मषहारिभिश्च}
{स्नानं सुतोयैः समुदाचर त्वं}
{मृत्युञ्जयानन्तगुणाभिराम}% ॥ ९ ॥

\twolineshloka
{आनीतेनातिशुभ्रेण कौशेयेनामरद्रुमात्}
{मार्जयामि जटाभारं शिव मृत्युञ्जय प्रभो}% ॥ १० ॥

\twolineshloka
{नानाहेमविचित्राणि चीरचीनाम्बराणि च}
{विविधानि च दिव्यानि मृत्युञ्जय सुधारय}% ॥ ११ ॥

\fourlineindentedshloka
{विशुद्धमुक्ताफलजालरम्यं}
{मनोहरं काञ्चनहेमसूत्रम्}
{यज्ञोपवीतं परमं पवित्रम्}
{आधत्स्व मृत्युञ्जय भक्तिगम्य}% ॥ १२ ॥

\fourlineindentedshloka
{श्रीगन्धं घनसारकुङ्कुमयुतं कस्तूरिकापूरितं}
{कालेयेन हिमाम्बुना विरचितं मन्दारसंवासितम्}
{दिव्यं देवमनोहरं मणिमये पात्रे समारोपितं}
{सर्वाङ्गेषु विलेपयामि सततं मृत्युञ्जय श्रीविभो}% ॥ १३ ॥

\twolineshloka
{अक्षतैर्धवलैर्दिव्यैः सम्यक्तिलसमन्वितैः}
{मृत्युञ्जय महादेव पूजयामि वृषध्वज}% ॥ १४ ॥

\fourlineindentedshloka
{चम्पकपङ्कजकुरवक-}
{कुन्दैः करवीरमल्लिकाकुसुमैः}
{विस्तारय निजमकुटं}
{मृत्युञ्जय पुण्डरीकनयनाप्त}% ॥ १५ ॥

\twolineshloka
{माणिक्यपादुकाद्वन्द्वे मौनिहृत्पद्ममन्दिरे}
{पादौ सत्पद्मसदृशौ मृत्युञ्जय निवेशय}% ॥ १६ ॥

\fourlineindentedshloka
{माणिक्यकेयूरकिरीटहारैः}
{काञ्चीमणिस्थापितकुण्डलैश्च}
{मञ्जीरमुख्याभरणैर्मनोज्ञैः}
{अङ्गानि मृत्युञ्जय भूषयामि}% ॥ १७ ॥

\fourlineindentedshloka
{गजवदनस्कन्दधृते-}
{नातिस्वच्छेन चामरयुगेन}
{गलदलकाननपद्मं}
{मृत्युञ्जय भावयामि हृत्पद्मे}% ॥ १८ ॥

\fourlineindentedshloka
{मुक्तातपत्रं शशिकोटिशुभ्रं}
{शुभप्रदं काञ्चनदण्डयुक्तम्}
{माणिक्यसंस्थापितहेमकुम्भं}
{सुरेश मृत्युञ्जय तेऽर्पयामि}% ॥ १९ ॥

\fourlineindentedshloka
{मणिमुकुरे निष्पटले}
{त्रिजगद्गाढान्धकारसप्ताश्वे}
{कन्दर्पकोटिसदृशं}
{मृत्युञ्जय पश्य वदनमात्मीयम्}% ॥ २० ॥

\fourlineindentedshloka
{कर्पूरचूर्णं कपिलाज्यपूतं}
{दास्यामि कालेयसमन्वितं च}
{समुद्भवं पावनगन्धधूपितं}
{मृत्युञ्जयाङ्गं परिकल्पयामि}% ॥ २१ ॥

\fourlineindentedshloka
{वर्तित्रयोपेतमखण्डदीप्त्या}
{तमोहरं बाह्यमथान्तरं च}
{साज्यं समस्तामरवर्गहृद्यं}
{सुरेश मृत्युञ्जय वंशदीपम्}% ॥ २२ ॥

\fourlineindentedshloka
{राजान्नं मधुरान्वितं च मृदुलं माणिक्यपात्रे स्थितं}
{हिङ्गूजीरकसन्मरीचिमिलितैः शाकैरनेकैः शुभैः}
{शाकं सम्यगपूपसूपसहितं सद्योघृतेनाप्लुतं}
{श्रीमृत्युञ्जय पार्वतीप्रिय विभो सापोशनं भुज्यताम्}% ॥ २३ ॥

\fourlineindentedshloka
{कूश्माण्डवार्ताकपटोलिकानां}
{फलानि रम्याणि च कारवल्ल्याः}
{सुपाकयुक्तानि ससौरभाणि}
{श्रीकण्ठ मृत्युञ्जय भक्षयेश}% ॥ २४ ॥

\twolineshloka
{शीतलं मधुरं स्वच्छं पावनं वासितं लघु}
{मध्ये स्वीकुरु पानीयं शिव मृत्युञ्जय प्रभो}% ॥ २५ ॥

\twolineshloka
{शर्करामिलितं स्निग्धं दुग्धान्नं गोघृतान्वितम्}
{कदलीफलसम्मिश्रं भुज्यतां मृत्युसंहर}% ॥ २६ ॥

\fourlineindentedshloka
{केवलमतिमाधुर्यं}
{दुग्धैः स्निग्धैश्च शर्करामिलितैः}
{एलामरीचमिलितं}
{मृत्युञ्जय देव भुङ्क्ष्व परमान्नम्}% ॥ २७ ॥

\fourlineindentedshloka
{रम्भाचूतकपित्थकण्टकफलैर्द्राक्षारसस्वादुमत्-}
{खर्जूरैर्मधुरेक्षुखण्डशकलैः सन्नारिकेलाम्बुभिः}
{कर्पूरेण सुवासितैर्गुडजलैर्माधुर्ययुक्तैर्विभो}
{श्रीमृत्युञ्जय पूरय त्रिभुवनाधारं विशालोदरम्}% ॥ २८ ॥

\fourlineindentedshloka
{मनोज्ञरम्भावनखण्डखण्डितान्}
{रुचिप्रदान् सर्षपजीरकांश्च}
{ससौरभान् सैन्धवसेवितांश्च}
{गृहाण मृत्युञ्जय लोकवन्द्य}% ॥ २९ ॥

\fourlineindentedshloka
{हिङ्गूजीरकसहितं}
{विमलामलकं कपित्थमतिमधुरम्}
{बिसखण्डाँल्लवणयुतान्}
{मृत्युञ्जय तेऽर्पयामि जगदीश}% ॥ ३० ॥

\fourlineindentedshloka
{एलाशुण्ठीसहितं}
{दध्यन्नं चारुहेमपात्रस्थम्}
{अमृतप्रतिनिधिमाढ्यं}
{मृत्युञ्जय भुज्यतां त्रिलोकेश}% ॥ ३१ ॥

\fourlineindentedshloka
{जम्बीरनीराञ्चितशृङ्गबेरं}
{मनोहरानम्लशलाटुखण्डान्}
{मृदूपदंशान् सहसोपभुङ्क्ष्व}
{मृत्युञ्जय श्रीकरुणासमुद्र}% ॥ ३२ ॥

\fourlineindentedshloka
{नागररामठयुक्तं}
{सुललितजम्बीरनीरसम्पूर्णम्}
{मथितं सैन्धवसहितं}
{पिब हर मृत्युञ्जय क्रतुध्वंसिन्}% ॥ ३३ ॥

\fourlineindentedshloka
{मन्दारहेमाम्बुजगन्धयुक्तैः}
{मन्दाकिनीनिर्मलपुण्यतोयैः}
{गृहाण मृत्युञ्जय पूर्णकाम}
{श्रीमत्परापोशनमभ्रकेश}% ॥ ३४ ॥

\fourlineindentedshloka
{गगनधुनीविमलजलैः}
{मृत्युञ्जय पद्मरागपात्रगतैः}
{मृगमदचन्दनपूर्णैः}
{प्रक्षालय चारु हस्तपदयुग्मम्}% ॥ ३५ ॥

\twolineshloka
{पुन्नागमल्लिकाकुन्दवासितैर्जाह्नवीजलैः}
{मृत्युञ्जय महादेव पुनराचमनं कुरु}% ॥ ३६ ॥

\fourlineindentedshloka
{मौक्तिकचूर्णसमेतैः}
{मृगमदघनसारवासितैः पूगैः}
{पर्णैः स्वर्णसमानैः}
{मृत्युञ्जय तेऽर्पयामि ताम्बूलम्}% ॥ ३७ ॥

\fourlineindentedshloka
{नीराजनं निर्मलदीप्तिमद्भिः}
{दीपाङ्कुरैरुज्ज्वलमुच्छ्रितैश्च}
{घण्टानिनादेन समर्पयामि}
{मृत्युञ्जयाय त्रिपुरान्तकाय}% ॥ ३८ ॥

\fourlineindentedshloka
{विरिञ्चिमुख्यामरवृन्दवन्दिते}
{सरोजमत्स्याङ्कितचक्रचिह्निते}
{ददामि मृत्युञ्जय पादपङ्कजे}
{फणीन्द्रभूषे पुनरर्घ्यमीश्वर}% ॥ ३९ ॥

\fourlineindentedshloka
{पुन्नागनीलोत्पलकुन्दजाती-}
{मन्दारमल्लीकरवीरपङ्कजैः}
{पुष्पाञ्जलिं बिल्वदलैस्तुलस्या}
{मृत्युञ्जयाङ्घ्रौ विनिवेशयामि}% ॥ ४० ॥

\fourlineindentedshloka
{पदे पदे सर्वतमोनिकृन्तनं}
{पदे पदे सर्वशुभप्रदायकम्}
{प्रदक्षिणं भक्तियुतेन चेतसा}
{करोमि मृत्युञ्जय रक्ष रक्ष माम्}% ॥ ४१ ॥

\fourlineindentedshloka
{नमो गौरीशाय स्फटिकधवलाङ्गाय च नमो}
{नमो लोकेशाय स्तुतविबुधलोकाय च नमः}
{नमः श्रीकण्ठाय क्षपितपुरदैत्याय च नमो}
{नमः फालाक्षाय स्मरमदविनाशाय च नमः}% ॥ ४२ ॥

\fourlineindentedshloka
{संसारे जनितापरोगसहिते तापत्रयाक्रन्दिते}
{नित्यं पुत्रकलत्रवित्तविलसत्पाशैर्निबद्धं दृढम्}
{गर्वान्धं बहुपापवर्गसहितं कारुण्यदृष्ट्या विभो}
{श्रीमृत्युञ्जय पार्वतीप्रिय सदा मां पाहि सर्वेश्वर}% ॥ ४३ ॥

\fourlineindentedshloka
{सौधे रत्नमये नवोत्पलदलाकीर्णे च तल्पान्तरे}
{कौशेयेन मनोहरेण धवलेनाच्छादिते सर्वशः}
{कर्पूराञ्चितदीपदीप्तिमिलिते रम्योपधानद्वये}
{पार्वत्याः करपद्मलालितपदं मृत्युञ्जयं भावये}% ॥ ४४ ॥

\fourlineindentedshloka
{चतुश्चत्वारिंशद्विलसदुपचारैरभिमतैः}
{मनःपद्मे भक्त्या बहिरपि च पूजां शुभकरीम्}
{करोति प्रत्यूषे निशि दिवसमध्येऽपि च पुमान्}
{प्रयाति श्रीमृत्युञ्जयपदमनेकाद्भुतपदम्}% ॥ ४५ ॥

\fourlineindentedshloka
{प्रातर्लिङ्गमुमापतेरहरहः सन्दर्शनात् स्वर्गदं}
{मध्याह्ने हयमेधतुल्यफलदं सायन्तने मोक्षदम्}
{भानोरस्तमये प्रदोषसमये पञ्चाक्षराराधनं}
{तत्कालत्रयतुल्यमिष्टफलदं सद्योऽनवद्यं दृढम्}% ॥ ४६ ॥

॥इति श्रीमत्परमहंसपरिव्राजकाचार्यस्य श्री-गोविन्द-भगवत्पूज्य-पाद-शिष्यस्य
श्रीमच्छङ्करभगवतः कृतौ मृत्युञ्जयमानसपूजास्तोत्रं सम्पूर्णम्॥
